\par\noindent\textbf{(1)} 机械能给 \(v=\sqrt{2gy}\)，线元 \(ds=\sqrt{1+y'^2}\,dx\)，故运动时间 \(t=(2g)^{-1/2}\int\sqrt{(1+y'^2)/y}\,dx\)。被积函数不显含 \(x\)，可用Beltrami第一积分，需对起点 \(y=0\) 的可积奇性单独处理。
\par\noindent\textbf{(2)} 第一积分 \(F-y'F_{y'}=[y(1+y'^2)]^{-1/2}=1/\sqrt{2a}\)，即 \(y(1+y'^2)=2a\)。令 \(y=a(1-\cos\theta)\)，则 \(dx/d\theta=a(1-\cos\theta)\)，积分得 \(x=a(\theta-\sin\theta)\)，为向下旋轮线。
\par\noindent\textbf{(3)} 从参数解计算 \(ds/d\theta=2a\sin(\theta/2)\)，能量速率 \(v=2\sqrt{ga}\sin(\theta/2)\)，故 \(dt/d\theta=\sqrt{a/g}\)。从 \(\theta=0\) 到 \(2\pi\) 的全拱运动时间 \(t=2\pi\sqrt{a/g}\)，不存在末端发散。
\par\noindent\textbf{(4)} 对旋轮线最低点定义有符号弧长 \(s=-4a\cos(\theta/2)\)，重力切向分量与 \(\sin(\theta/2)\) 成正比，严格得到 \(\ddot s+[g/(4a)]s=0\)。因此在同一拱内往返最低点的周期 \(4\pi\sqrt{a/g}\) 与幅度无关；轨道和初始条件必须符合无摩擦约束。
